\section{Introduction}\label{sec:intro}

A conservative reading of the causal-set and order-first programs is that,
for future- and past-distinguishing Lorentzian spacetimes, causal or
chronological structure determines the conformal class. Supplying a volume
element fixes the remaining conformal factor, up to diffeomorphism and unit
conventions \cite{blms1987,malament1977,hkm1976,kronheimer1967,sorkin2005,surya2019}.
This paper treats that statement operationally and quantitatively. Rather
than ask whether spacetime ``is'' a causal set, we ask: fix a causal (or
null-accessibility) order in a known model; then, supplying an explicitly
declared ingredient set, which spacetime quantities can a finite procedure
actually reconstruct, with what error, and where does each reconstruction
stop?

The value of posing it this way is an explicit \emph{ledger}. Many
individual reconstructions here are standard---proper time from longest
chains \cite{brightwell1991}, volume from interval cardinality
\cite{blms1987}, dimension from ordering fractions \cite{myrheim1978,meyer1988},
radar coordinates from an observer protocol. What is easy to lose, and what
we make explicit, is exactly which extra structure each requires and what
remains underdetermined without it. The ladder also isolates three
instructive negative results: causal order alone does not fix conformal
scale; a single observer yields only unsigned distance (a reflection
degeneracy); and finite signal speed alone does not imply Lorentzian
structure.

The contributions are: (i) a reconstruction ladder that indexes recoverable
spacetime quantities by the minimal supplied ingredient set, with per-rung
error behaviour in controlled models; (ii) grounded validations of dimension,
proper time, radar decomposition, Lorentz-map and atlas consistency, and
measure-dependent volume; (iii) three bounding negative results (conformal
scale, reflection degeneracy, finite-speed lattice); (iv) a Rindler
reconstruction-inaccessibility analogue; and (v) a preregistered
$(3{+}1)$D capstone validation that a pure-Weyl deformation of the light
cones---at identical coordinate volume form and sprinkling law---measurably
moves the causal order at finite density (\sref{sec:capstone}), together
with its type-D (Schwarzschild) extension, where separate preregistered
stages confirm the paired claim and detect single-poset discrimination
(\sref{sec:schwarzschild}); and (vi) a certified diamond-volume enclosure
and a preregistered, prediction-anchored Poisson-count validation across a
four-rung Schwarzschild mass ladder (\sref{sec:count}), with an auxiliary
instrument audit of the oracle stack (\sref{sec:oracle}). All results are
controlled validations, not evidence that geometry reduces to order.

A literature priority search (August 2026) found partial overlap in four
areas: causal-order-and-number/volume reconstruction (Braun
2025~\cite{braun2025}; Myrheim 1978~\cite{myrheim1978}; the number-volume
correspondence itself is Sorkin's~\cite{sorkin1997forks}), Schwarzschild
sprinkle-count verification (Hom\v{s}ak--Veroni~\cite{homsakveroni2024}
\S V.1), analytic diamond-volume results
(Berthiere--Gibbons--Solodukhin~\cite{berthiere2016};
Roy--Sinha--Surya~\cite{roysinhasurya2013}), and Schwarzschild
causal-relation algorithms (He--Rideout~\cite{herideout2009}; for
sprinkling-and-relations tooling more broadly,
Cunningham--Krioukov~\cite{cunninghamkrioukov2018}). The directed-rounding
volume enclosure and the preregistered equivalence gate against it appear to
have no direct precedent in that search (\sref{sec:count}); the parallel
Alfyorov--Shnyukov
estimator~\cite{alfyorov2026weyl} measures Weyl curvature from Hasse
diagrams and is taken up alongside the capstone's scope in
\sref{sec:notestablish}.
